\section{Asymmetric Hermite nodes have the same precision optimum}
\label{sec:asymmetric-hermite}

The earlier exact minimax theorem treats valuation-symmetric
multiplicities~\cite{precision}. An earlier barrier explicitly leaves
asymmetric node selection open~\cite{analysis}. We remove that restriction
for the \emph{fixed Hermite reconstruction with independent jet errors}.
This is a precision-certificate theorem. It does not assert that the
jets attainable by the source deformation can vary independently, or
rule out recovery that exploits their correlations.

\subsection{Exact loss of a node's recovery block}

Let $a\ge u\ge1$, $N=a+1$, $L=2^a$, and let $S$ be a complete
integer residue system modulo $2^N$ containing zero. For example,
$S=\{-L,\ldots,L-1\}$. Assign nonnegative integer multiplicities
$q_t$ to the nonzero nodes, put $q_0=0$, and require $q_t=0$ when
$v_2(t)\ge u$. There is no symmetry or monotonicity requirement.

For an active node $t$, put $i=v_2(t)$, $D_t=2^{a-i}$ and
\[
 E_t(T)=\prod_{v\in S\setminus\{0,t\}}
       \left(\frac{T-v}{t-v}\right)^{q_v}.
\]
The ordinary Hermite reconstruction of $F(0)$ uses the first $q_t$
coefficients of $F(t+2^Nz)$. The coefficient of its highest raw jet is
exactly
\[
 E_t(0)(-t/2^N)^{q_t-1}.
\]
Indeed, $E_t(t+2^Nz)^{-1}$ is an integral power series with constant
term one. The top raw jet only contributes through that constant term.
The negative valuation of the displayed coefficient is
\begin{equation}\label{eq:lambda-node}
 \Lambda_t(q)=
 \sum_{v\ne0,t}q_v\bigl(v_2(t-v)-v_2(v)\bigr)
 +(q_t-1)(N-i).
\end{equation}
Every other raw-jet coefficient is a sum of terms
$E_t(0)c_j(-t/2^N)^h$ with $c_j$ integral and $h\le q_t-1$.
Its loss is therefore at most $\Lambda_t(q)$. The highest jet attains
this bound exactly. Consequently
\[
 p+\max\{0,\Lambda_t(q)\}
\]
is the exact nondecreasing precision allowance for this node's
independent-error recovery block at output precision $p$.
Define its normalized worst loss by
\begin{equation}\label{eq:H-objective}
 H(q)=\max\left(0,\max_{t:q_t>0}\frac{\Lambda_t(q)}{D_t}\right).
\end{equation}

For $1\le k\le N$, define the residue-subtree mass
\[
 S_k(t)=\sum_{\substack{v\in S\\v\equiv t\pmod{2^k}}}q_v.
\]
Expressing valuations as sums of congruence indicators gives
\begin{equation}\label{eq:subtree-loss}
 \Lambda_t(q)=
 \sum_{k=i+1}^{N}\bigl(S_k(t)-S_k(0)\bigr)-(N-i).
\end{equation}
To verify the endpoint carefully, for $k<N$ excluding $v=t$
subtracts $(N-1-i)q_t$ from the valuation-difference sum.
The highest-jet term adds $(q_t-1)(N-i)$, leaving
$q_t-(N-i)$. The $q_t$ is precisely the $k=N$ subtree term.

\subsection{Inactive nodes and convexity}

Equation~\eqref{eq:lambda-node} can be evaluated as an affine expression
even at inactive nodes. We also allow nonnegative real $q_t$ as a
relaxation; real multiplicities are not actual jet counts.

\begin{lemma}\label{lem:inactive}
For every such nonnegative vector,
\[
 H(q)=
 \max\left(0,\max_{\substack{t\in S\setminus\{0\}\\v_2(t)<u}}
                    \frac{\Lambda_t(q)}{D_t}\right).
\]
In particular the relaxed objective is convex.
\end{lemma}
\begin{proof}
Within an annulus $v_2(t)=i$, both $D_t$ and the negative terms in
\eqref{eq:subtree-loss} are constant. Its maximization therefore depends
only on the positive path sum $\sum_{k=i+1}^NS_k(t)$.
If an inactive leaf has an ancestor containing positive mass, take
the smallest such ancestor and an active leaf inside it. The two
paths agree above that ancestor. Below it the inactive path has
zero mass and the active path eventually has positive mass.
The active leaf has a strictly larger path sum.
Thus a nonempty annulus attains its maximum at an active leaf.
An empty annulus has zero positive path mass and strictly negative
affine loss. Adjoining zero makes its contribution irrelevant.
The displayed expression is a maximum of affine functions.
\end{proof}

\begin{theorem}[Symmetrization of arbitrary node multiplicities]
\label{thm:hermite-symmetry}
Replace every $q_t$ by the average multiplicity on its valuation
annulus, obtaining $\overline q$. Then
\[
 \sum_t\overline q_t=\sum_tq_t,\qquad
 H(\overline q)\le H(q).
\]
The real minimum of this precision certificate is therefore attained
by a valuation-symmetric profile, even when arbitrary asymmetric
node choices are allowed.
\end{theorem}
\begin{proof}
Consider automorphisms of the complete binary residue tree that fix
the leaf zero. They preserve $v_2(t)$ and $v_2(t-v)$, and act
transitively on every annulus. Permuting the multiplicity vector by
any such automorphism preserves \eqref{eq:H-objective} and total mass.
Averaging over this finite group gives $\overline q$. Convexity from
Lemma~\ref{lem:inactive} proves the inequality.
Only the valuations of the Hermite coefficients are used. We do not
assume that the source polynomial deformation itself is invariant
under these tree automorphisms.
\end{proof}

\subsection{An exact finite minimax formula}

For a symmetric profile, write $q_i$ for its common value at
valuation $i$ and set $q_i=0$ for $i\ge u$. Let
$D_i=2^{a-i}$ and $n_i=N-i$. The exact normalized loss is
\begin{equation}\label{eq:finite-symmetric}
 \frac{\Lambda_i}{D_i}
 =\left(2-\frac1{D_i}\right)q_i
  -\sum_{j=1}^{u-1-i}\frac{j}{2^j}q_{i+j}
  -\frac{n_i}{D_i}.
\end{equation}
For instance, the finite normalization sum has diagonal
$2-(n_i+1)/D_i$; adding $(q_i-1)n_i/D_i$ gives precisely
\eqref{eq:finite-symmetric}.

Given required normalized total mass $W>0$, define backward rational
recurrences
\begin{align}
 r_i&=\frac{1+\sum_{j=1}^{u-1-i}j2^{-j}r_{i+j}}
                   {2-D_i^{-1}},\label{eq:r-rec}\\
 c_i&=\frac{n_i/D_i+\sum_{j=1}^{u-1-i}j2^{-j}c_{i+j}}
                   {2-D_i^{-1}},\label{eq:c-rec}\\
 R_*&=\sum_{i=0}^{u-1}2^{-i}r_i,\qquad
 C_*=\sum_{i=0}^{u-1}2^{-i}c_i.\label{eq:RC}
\end{align}
The starred $C_*$ here is a local minimax quantity, unrelated to
the algorithmic constants $C_T,C_S$.

\begin{theorem}[Finite optimum and integer enclosure]
\label{thm:finite-hermite}
Over all nonnegative real node multiplicities supported on
$v_2(t)<u$ and satisfying $\sum_tq_t\ge LW$,
\[
 \min H(q)=\eta_*=\max\left\{0,\frac{W-C_*}{R_*}\right\}.
\]
If $W\ge C_*$, a minimizing symmetric profile is
$q_i=\eta_*r_i+c_i$. If $W<C_*$, use $q_i=(W/C_*)c_i$.
The optimum over integer multiplicities, even allowing arbitrary
node-by-node asymmetry, lies in $[\eta_*,\eta_*+2)$.
\end{theorem}
\begin{proof}
Theorem~\ref{thm:hermite-symmetry} reduces the real lower bound to
symmetric profiles. If $H(q)\le\eta$, equation
\eqref{eq:finite-symmetric} gives
\[
 q_i\le
 \frac{\eta+n_i/D_i+\sum_{j=1}^{u-1-i}j2^{-j}q_{i+j}}
      {2-D_i^{-1}}.
\]
Backward induction yields $q_i\le\eta r_i+c_i$. Since an annulus
contains $D_i=L2^{-i}$ nodes, completeness requires
$W\le\sum_i2^{-i}q_i\le\eta R_*+C_*$. This proves the stated
lower bound. The displayed profile achieves equality in every
constraint if $W\ge C_*$. Scaling $c_i$ downward when $W<C_*$
makes every affine loss nonpositive and gives exact mass $W$.

For the integer upper bound, round every component of the chosen
symmetric real profile upward. The total mass cannot decrease.
In \eqref{eq:finite-symmetric}, the own-coordinate increase is
less than $2-D_i^{-1}<2$, while every other increase has a
nonpositive coefficient. Thus $H$ rises by less than two.
The unrestricted integer minimum is no smaller than its real
relaxation, which proves both sides.
\end{proof}

\subsection{The leading constant six survives asymmetry}

Let $D_{\min}=2^{a-u+1}$. In the limit $D_i=\infty$, the backward
recurrence has the explicit solution, already implicit in the
symmetric taper analysis~\cite{precision},
\[
 r_i^{(\infty)}
 =\frac{u-i}{6}+\frac5{18}+\frac29\,4^{-(u-i)},\qquad
 R_\infty=\frac u3+\frac29-\frac29\,4^{-u}.
\]
Positivity of the finite recurrences gives
\begin{equation}\label{eq:R-comparison}
 R_\infty\le R_*
 \le\left(1-\frac1{2D_{\min}}\right)^{-u}R_\infty,
 \qquad
 0\le C_*\le\frac{a+1}{D_{\min}}R_*.
\end{equation}
For the first estimate, expand the backward recurrence as a sum over
increasing index paths. Each visited index changes the limiting
contribution by a factor at most
$(1-1/(2D_{\min}))^{-1}$, and a path visits at most $u$ indices.
The second estimate follows by comparing the inhomogeneous terms
in \eqref{eq:r-rec} and \eqref{eq:c-rec}.

Hermite recovery for degree at most $M(2L-1)$ requires total
multiplicity at least $M(2L-1)+1$, so
\[
 W=\frac{M(2L-1)+1}{L}=(2+o(1))M.
\]
If $u\to\infty$, $u/D_{\min}\to0$,
$(a+1)/D_{\min}=o(M/u)$, and $M/u\to\infty$, equations
\eqref{eq:R-comparison} and Theorem~\ref{thm:finite-hermite} imply
\begin{equation}\label{eq:asymmetric-six}
 \min_{\text{integer node profiles}}H(q)
   =(6+o(1))\frac M u.
\end{equation}
The thin-gap schedules have $D_{\min}\ge\lambda^6$ and
$a,u=O(\lambda)$, so these conditions hold at active calls.
Arbitrary omission or asymmetric weighting of nodes within the
retained valuation levels therefore cannot improve the leading
$MD/u$ precision scale in this fixed independent-error certificate.

The exact finite checker tests the valuation identity, the inactive-node
extension, symmetrization, and the minimax construction. Its role is
to detect finite algebraic mistakes; the general theorem is proved
above rather than inferred from those checks.
